% Silicon Cartography: Constraint-First Programming for Spatial Chips
% Author: Flyxion
% Typeset with LuaLaTeX / XeLaTeX

\documentclass[11pt, a4paper]{article}

\usepackage{fontspec}
\setmainfont{TeX Gyre Pagella}
\usepackage{microtype}
\usepackage{geometry}
\geometry{a4paper, top=3cm, bottom=3cm, left=3.2cm, right=3.2cm}

\usepackage{amsmath, amssymb, amsthm}
\usepackage{mathtools}
\usepackage{hyperref}
\hypersetup{
  colorlinks=true,
  linkcolor=black,
  citecolor=black,
  urlcolor=black,
  pdftitle={Silicon Cartography\\
\large Constraint-First Programming for Spatial Chips},
  pdfauthor={Flyxion}
}
\usepackage{setspace}
\setstretch{1.15}
\usepackage{booktabs}
\usepackage{enumitem}
\setlist[itemize]{noitemsep, topsep=4pt}
\usepackage{titlesec}
\titleformat{\section}{\large\bfseries}{\thesection.}{0.8em}{}
\titleformat{\subsection}{\normalsize\bfseries}{\thesubsection.}{0.8em}{}
\usepackage{abstract}
\renewcommand{\abstractnamefont}{\normalfont\bfseries}
\usepackage{parskip}
\setlength{\parskip}{0.6em}
\setlength{\parindent}{0pt}

\newtheorem{definition}{Definition}
\newtheorem{theorem}{Theorem}
\newtheorem{proposition}{Proposition}
\newtheorem{remark}{Remark}

\title{\textbf{Silicon Cartography}\\
\large Constraint-First Programming for Spatial Chips}
\author{Flyxion\\[0.3em]
\small Independent Researcher}
\date{May 2026}

% -----------------------------------------------------------
\begin{document}
\maketitle
\thispagestyle{empty}

% -----------------------------------------------------------
\begin{abstract}
Contemporary compilation assumes a computational substrate that is abundant, uniform,
and energetically indifferent. This assumption fails for a growing class of spatial chips:
multi-core mesh architectures, neuromorphic arrays, FPGAs, and low-power embedded systems
in which energy, memory locality, routing topology, and inter-node communication are the
dominant engineering parameters. We introduce \emph{Silicon Cartography}, a
constraint-first programming discipline and corresponding compiler architecture that
begins not with a program but with an interrogation of the machine. The hardware is modeled
as a constrained spatial computer---a metrized graph $\mathcal{H} = (N, E, \mathcal{A},
\mathcal{M}, \mathcal{C}, \Phi)$ encoding nodes, communication edges, instruction algebras,
memory geometry, hard constraints, and physical cost functionals. A program is then a
mapping problem: factor a desired computation into locally self-contained kernels, place
those kernels onto the discovered hardware geometry, route communication, and schedule
wake/sleep behavior so that only necessary physical transitions occur.

We develop this framework across five interrelated dimensions. First, we establish the
formal structure of the hardware graph and its compilation functor, including feasibility,
optimality, and compositionality results. Second, we develop a taxonomy of fractal
procedural chip topologies---fractal island meshes, tree-of-meshes, Hilbert-curve
chips, Sierpinski compute fabrics, dragon-curve pipelines, and Apollonian heterogeneous
fabrics---each with characteristic locality properties and programming disciplines.
Third, we analyze the communication manifold structure of spatial chips through
graph Laplacians, spectral connectivity, and synchronization sheaves, connecting deadlock
detection to Čech cohomology obstruction classes. Fourth, we introduce a procedural
topology language that separates hardware description from compilation policy, enabling
systematic design-space exploration. Fifth, we develop a visualization framework that
represents the six interacting dimensions of spatial chip behavior---topology, activity,
energy, communication, memory, and time---through activity heatmaps, event wave
propagation, causal braid diagrams, multiscale semantic zooming, and energy landscape
terrain maps. Finally, we ground the energy-minimality criterion in thermodynamic
computation theory, connecting Moore's principle that every unnecessary transistor
transition is waste to Landauer's fundamental bound on the thermodynamic cost of
logically irreversible operations.
\end{abstract}

\newpage

% -----------------------------------------------------------
\tableofcontents
\newpage

% -----------------------------------------------------------
\section{Motivation: The Return of Physical Constraints}

The dominant tradition in software engineering treats hardware as a uniform abstraction.
A program is written against an idealized machine---sequential, byte-addressable,
with abundant RAM and an optimizing compiler mediating between intent and execution.
This tradition emerged from a historically contingent fact: from roughly 1970 to 2010,
transistor density, clock speed, and memory bandwidth improved rapidly enough that
programmers could afford to ignore the physical substrate. Abstraction layers---operating
systems, virtual machines, garbage collectors, high-level languages---accumulated precisely
because the hardware could absorb their overhead.

That era is structurally over in several important engineering domains. Edge inference,
wearable computing, embedded intelligence, sensor swarms, and large-scale datacenter
deployments each face a common bottleneck: energy. As Chuck Moore observed in his
documentation of the GreenArrays architecture, a 144-node mesh of tiny Forth computers,
power consumption is no longer a secondary engineering parameter but the primary one.
Even two logically equivalent delay loops may consume measurably different energy because
one toggles more decode circuitry than the other. At the spatial and temporal scales of
modern embedded systems, every unnecessary transistor transition carries a thermodynamic cost.

The implications for programming methodology are radical. If energy is primary, then the
compiler's job is not to translate syntax while optimizing for speed. It is to find a
physical placement of computation that minimizes unnecessary state transitions, routes
communication along low-cost paths, keeps idle hardware in low-energy dormancy, and
exploits the specific geometry of the chip's memory, bus, and interconnect topology.
This requires the compiler---and the programmer---to begin by knowing the machine.

Silicon Cartography is the discipline of beginning there: with an interrogation of the
hardware geometry before a single line of code is written.

% -----------------------------------------------------------
\section{The Hardware Model}

\subsection{Spatial Chips as Constraint Graphs}

We define a \emph{spatial chip} as any computational substrate in which the topology
of processing units, their communication paths, and their local memory resources are
physically fixed and non-uniform. GreenArrays-style mesh processors, FPGAs, neuromorphic
arrays, and many modern embedded SoCs satisfy this description. The key property is that
computation has \emph{location}: a kernel running on node $n_i$ has access to different
resources at different costs than the same kernel running on node $n_j$, and communication
between distant nodes is more expensive than communication between neighbors.

\begin{definition}[Hardware Graph]
A \emph{hardware graph} is a tuple
\[
H = (N,\, E,\, I,\, M,\, C,\, P)
\]
where:
\begin{itemize}
  \item $N$ is a finite set of \emph{computational nodes} (cores, processing elements);
  \item $E \subseteq N \times N$ is a set of directed \emph{communication edges}
        representing physical links, buses, DMA channels, or shared memory paths;
  \item $I : N \to \mathcal{P}(\text{Instr})$ assigns to each node a set of executable
        instructions or operations;
  \item $M : N \to \mathbb{N}$ assigns to each node a local memory capacity (words or bytes);
  \item $C$ is a set of \emph{hard constraints}: maximum fanout, routing exclusions,
        timing windows, voltage domains, synchronization requirements; and
  \item $P : N \cup E \times \text{Op} \to \mathbb{R}_{\geq 0}$ is a \emph{physical cost model}
        assigning energy, latency, heat, and contention costs to operations and
        communication events.
\end{itemize}
\end{definition}

The hardware graph is not a datasheet. It is an \emph{executable constraint map}: a
structure that can be queried to answer questions of the form ``can computation $k$ run
at node $n$?'', ``what does it cost to send message $m$ along edge $e$?'', and
``which nodes can sleep while kernel $k'$ executes?''. The first task of the Silicon
Cartographer system is to construct this map.

\subsection{The Cartographer Layer}

Construction of $H$ may proceed by several means depending on the chip and available
tooling. For chips with published hardware description languages (VHDL, SystemVerilog,
device tree specifications), the cartographer parses these to extract node counts,
interconnect topology, instruction sets, and power domains. For underdocumented or novel
chips, it may issue low-level probe sequences---latency measurements, instruction timing
loops, memory bandwidth tests, power consumption sampling via external instrumentation---to
empirically derive the cost model $P$.

The result in all cases is a machine-readable description of the chip's \emph{physical
possibility space}: what computations are locally admissible at each node, what
communication patterns are feasible across each edge, and what physical costs attach
to each operation. This is the foundation on which all subsequent compilation rests.

\subsection{Node States and Energy Topology}

An important dimension of the hardware model that static datasheets tend to omit is
the \emph{energy topology}: the graph of power states that each node may occupy and
the transition costs between them. For a GreenArrays node, this includes at minimum a
full-execution state, a port-blocked waiting state (listening for a message from a neighbor),
and a deep-idle state. The waiting state consumes dramatically less energy than execution
but more than deep idle. Programming the chip well means keeping most nodes in low-energy
states most of the time, with brief excursions into execution triggered by communication
events.

We formalize this as a finite automaton $\mathcal{A}_n = (Q_n, \Sigma_n, \delta_n, q_0^n)$
for each node $n$, where $Q_n$ is the set of power states, $\Sigma_n$ is the set of
triggering events (instruction issue, message arrival, clock interrupt), $\delta_n$ is
the transition function, and $q_0^n$ is the default idle state. The cost model $P$ then
includes the energy cost of each transition $\delta_n(q, \sigma) = q'$ as
$P_{\text{trans}}(n, q, \sigma)$.

Minimizing total energy expenditure over a computation trace $\tau$ becomes:
\[
\text{minimize} \quad \sum_{n \in N} \sum_{t} P_{\text{trans}}(n,\, q_n(t),\, \sigma_n(t))
+ P_{\text{exec}}(n,\, \text{instr}_n(t))
+ \sum_{e \in E} P_{\text{comm}}(e,\, \text{msg}_e(t))
\]
subject to the constraint that the computation trace produces the correct output.

% -----------------------------------------------------------
\section{The Program Model}

\subsection{Kernels, Dependencies, and Requirements}

Given a high-level task specification, the program decomposer factors it into a
\emph{program graph}:

\begin{definition}[Program Graph]
A \emph{program graph} is a tuple
\[
G = (K,\, D,\, R)
\]
where:
\begin{itemize}
  \item $K$ is a finite set of \emph{kernels}---self-contained computational units
        that operate on local data, produce outputs, and communicate through well-defined
        message interfaces;
  \item $D \subseteq K \times K$ is a directed graph of \emph{data dependencies}: an
        edge $(k_i, k_j) \in D$ indicates that $k_j$ requires output from $k_i$; and
  \item $R : K \to \mathcal{C}$ assigns to each kernel its \emph{resource requirements}:
        instruction types needed, minimum local memory, required timing windows, accuracy
        or reliability constraints.
\end{itemize}
\end{definition}

Crucially, kernels are designed to be \emph{locally self-contained}. Each kernel
communicates only through its message interface---it does not assume a shared global
address space. This discipline directly mirrors the computational model of spatial chips,
where no global memory exists and all inter-node communication is explicit.

\subsection{The Decomposer}

The decomposer's task is to factor a task specification into $G$ in a way that respects
the chip's physical geometry. This requires answering a set of questions that are typically
invisible in mainstream programming:

\begin{itemize}
  \item What data can remain \emph{entirely local} to a single node, never appearing on
        any communication link?
  \item What data \emph{must move}, and along which communication paths?
  \item Which kernels can remain dormant until triggered by a message, and what is the
        expected message rate?
  \item What can be \emph{preloaded} onto a node's stack or local memory before
        computation begins?
  \item What can be \emph{delegated to a neighbor} acting as a smart memory or
        preprocessor, reducing the computation that the primary node must perform?
  \item What computation can be \emph{avoided entirely} by exploiting structure in the
        data (sparsity, symmetry, known ranges)?
\end{itemize}

These questions are not optimizations applied after a naive decomposition. They are
the primary structure of the decomposition itself. On a spatial chip, a design that
requires dense global communication is not merely inefficient---it may be physically
infeasible.

% -----------------------------------------------------------
\section{The Compilation Functor}

\subsection{Placement and Routing}

The compilation step finds a mapping $\varphi : G \to H$ that places kernels onto nodes,
routes dependencies onto communication edges, and assigns wake/sleep schedules consistent
with the chip's energy topology.

\begin{definition}[Compilation Mapping]
A \emph{valid compilation} of $G$ into $H$ is a pair $(\varphi_K, \varphi_D)$ where:
\begin{itemize}
  \item $\varphi_K : K \to N$ maps each kernel to a node such that
        $R(k) \subseteq I(\varphi_K(k))$ (the node can execute the kernel) and
        the memory required by $k$ does not exceed $M(\varphi_K(k))$;
  \item $\varphi_D : D \to \text{Paths}(E)$ maps each dependency $(k_i, k_j)$ to
        a directed path in $E$ from $\varphi_K(k_i)$ to $\varphi_K(k_j)$;
  \item for each edge $e \in E$, the aggregate communication load from all paths
        $\varphi_D(d)$ passing through $e$ does not exceed $e$'s bandwidth constraint; and
  \item all timing and synchronization constraints in $C$ are satisfied by the induced
        schedule.
\end{itemize}
\end{definition}

The cost of a compilation is:
\[
\text{Cost}(\varphi) = \alpha \cdot E_{\text{exec}}(\varphi)
+ \beta \cdot E_{\text{comm}}(\varphi)
+ \gamma \cdot L_{\text{critical}}(\varphi)
+ \delta \cdot M_{\text{pressure}}(\varphi)
+ \varepsilon \cdot S_{\text{waste}}(\varphi)
\]
where $E_{\text{exec}}$ is total execution energy, $E_{\text{comm}}$ is total
communication energy, $L_{\text{critical}}$ is critical-path latency,
$M_{\text{pressure}}$ is peak memory pressure, $S_{\text{waste}}$ is synchronization
waste (time spent in active waiting rather than dormancy), and $\alpha, \beta, \gamma,
\delta, \varepsilon \geq 0$ are domain-specific weighting coefficients.

For Chuck Moore's computational philosophy, the dominant constraint is $\alpha + \beta \gg
\gamma$: energy dominates over speed. For hard real-time systems, $\gamma$ dominates.
For memory-constrained microcontrollers, $\delta$ dominates. The cost structure is
chip-class-specific, and the cartographer derives appropriate defaults from the
hardware model.

\subsection{Code Generation}

Once a valid compilation $\varphi$ has been found, code generation is chip-specific
but structurally uniform. For each node $n \in N$, the code generator emits a program
that:

\begin{enumerate}
  \item preloads any required data onto $n$'s local stack or memory;
  \item enters a low-energy waiting state, blocked on the first expected input message;
  \item upon message arrival, executes the assigned kernel and produces output messages;
  \item routes output messages to $\varphi_D$-specified neighbors; and
  \item returns to the waiting state.
\end{enumerate}

For GreenArrays-style chips, this generates colorForth or Forth primitives that directly
encode the stack operations and port reads/writes. For FPGAs, it generates dataflow graphs
or HDL modules. For GPUs, it generates CUDA or OpenCL kernels with appropriate memory
hierarchy awareness. For neuromorphic chips, it generates spike-routing configurations.
The \emph{shape} of the generated code varies; the \emph{structure} of the generation
process is common.

% -----------------------------------------------------------
\section{Admissibility Geometry and RSVP Correspondence}

\subsection{The Chip as an Admissibility Field}

The hardware graph $H$ defines not merely a set of constraints but a \emph{topology of
reachability}: the set of computations that are physically possible, at what cost, and
along what paths. In the language of the RSVP (Relativistic Scalar-Vector Plenum)
framework, $H$ defines an \emph{admissibility field} $\mathcal{A}_H$ over the space
of computation-placements.

A computation-placement is a pair $(G, \varphi)$: a program graph together with a
specific mapping into the hardware. The admissibility field assigns to each such pair
a value in $[0, 1]$ reflecting whether the placement is feasible within the chip's
hard constraints, with the value decreasing as constraints are approached and reaching
zero outside the feasible region.

\begin{definition}[Admissibility Field]
The admissibility field of a hardware graph $H$ is a function
\[
\mathcal{A}_H : \{(G, \varphi)\} \to [0, 1]
\]
defined by:
\[
\mathcal{A}_H(G, \varphi) = \prod_{c \in C} \sigma_c(G, \varphi)
\]
where $\sigma_c$ is a soft satisfaction measure for constraint $c$, equal to $1$ when
$c$ is fully satisfied and $0$ when violated. Hard constraints enter as $\{0,1\}$-valued
factors; soft constraints (energy budgets, latency preferences) enter as smoothly
interpolated values.
\end{definition}

Compilation is then constrained projection: finding the highest-quality placement within
the admissible region,
\[
\varphi^* = \arg\max_{\varphi : \mathcal{A}_H(G,\varphi) = 1} \bigl(-\text{Cost}(\varphi)\bigr).
\]

\subsection{Instruction Execution as Local Entropy Expenditure}

The RSVP framework treats entropy not as a global thermodynamic quantity but as a
local, field-valued measure of irreversible state transitions. Instruction execution
on a spatial chip is a direct physical instantiation of this idea: each instruction
toggles a set of transistors, dissipating energy proportional to the number of
$0 \to 1$ and $1 \to 0$ transitions in the affected registers, buses, and decode
circuitry. The chip's physical cost model $P$ is essentially a local entropy expenditure
map.

Idle states correspond to low-transition basins: the node is in a metastable
configuration that changes slowly or not at all, minimizing entropy production.
Communication events are brief, spatially localized entropy expenditures that propagate
along admissibility channels (the edges $E$) to trigger state changes in neighboring
nodes.

Stack locality in Forth-style architectures corresponds to compressed trajectory memory:
the computation's recent history is encoded in the top few cells of the stack, which
are physically implemented as fast registers, minimizing the energy cost of accessing
state from the recent past.

\subsection{Blocking Behavior as Sparse Event Activation}

The blocking behavior that Moore describes---a node enters a low-energy waiting state
until a message arrives from a neighbor, then briefly executes and returns to
dormancy---is structurally identical to sparse-event cognition models in neuromorphic
computing and asynchronous neural signaling. In both cases, the energy cost of a
computation is proportional to the rate and magnitude of state-change events, not
to the passage of time per se.

In RSVP terms, the computation trace $\tau$ is a sequence of brief, spatially localized
trajectory segments separated by intervals of near-zero field activity. The chip is
not a continuously executing machine; it is a field of cooperative local operators that
activate opportunistically when boundary conditions (incoming messages, clock edges,
sensor events) make activation necessary.

\begin{proposition}[Energy Minimality as Trajectory Sparsity]
A computation $\tau$ on hardware graph $H$ is energy-minimal if and only if the
induced activation sequence $\mathcal{E}(\tau) = \{(n, t) : n \text{ executes at time } t\}$
is sparse in both the node dimension (few nodes active at any moment) and the temporal
dimension (each active node's active intervals are brief relative to its dormant intervals).
\end{proposition}

This is not merely a performance observation. It is a statement about the physical
structure of good programs on spatial chips: a well-written program on a GreenArrays-like
machine looks like a sparse activation field, not a dense computation trace.

% -----------------------------------------------------------
\section{Compiler Architecture}

\subsection{The Four Layers}

The Silicon Cartography compiler is organized into four layers, each with a well-defined
input/output contract.

\textbf{Layer 1: Hardware Cartographer.} Accepts chip description (HDL, device tree,
empirical probe results) and produces the hardware graph $H = (N, E, I, M, C, P)$.
Responsibilities: topology extraction, instruction set parsing, memory map construction,
power-state automaton derivation, cost model calibration.

\textbf{Layer 2: Capability Model.} Refines $H$ into an executable constraint map:
a queryable data structure supporting questions about local admissibility, edge capacity,
power-state transitions, and routing feasibility. This layer is the foundation of all
subsequent optimization; it must be consistent and queryable in logarithmic time.

\textbf{Layer 3: Program Decomposer.} Accepts a high-level task specification and the
capability model, and produces a program graph $G = (K, D, R)$ with kernel
interfaces designed to fit the chip's communication model. The decomposer applies
locality analysis (what can stay local?), communication analysis (what must move and
how often?), and sparsity analysis (what computation can be avoided?).

\textbf{Layer 4: Placement, Routing, and Code Generator.} Solves the constrained
optimization problem to find $\varphi^* : G \to H$, then emits chip-specific code
for each assigned node. The optimization may use constraint programming, simulated
annealing, or reinforcement learning depending on chip complexity; the code generator
is chip-class-specific.

\subsection{The Compiler's Role}

Chuck Moore's insistence that the compiler ``merely translates'' and performs no
optimization is not a limitation of Forth but a design philosophy: the programmer
should understand the hardware geometry directly, not delegate that understanding
to an optimizer. Silicon Cartography is consistent with this philosophy at the
strategic level while acknowledging that on complex multi-node chips, the placement
and routing problem is too large for unaided human cognition.

The resolution is that the compiler's optimization is \emph{structural}, not
\emph{semantic}. It does not rewrite the kernel logic; it finds the best physical
placement of kernels that the programmer has already designed to be hardware-aware.
The programmer still reasons about locality, communication, and energy. The compiler
automates the combinatorial search over physical placements.

In this sense, the Silicon Cartographer is less like a traditional optimizing compiler
and more like a \emph{constraint solver over hardware geometry}: it finds the mapping
$\varphi$ that satisfies all hard constraints and minimizes the energy cost model,
but it trusts the programmer to have designed kernels that are worth placing well.

% -----------------------------------------------------------
\section{Target Chip Classes}

The Silicon Cartography framework is designed to generalize across chip classes,
with chip-specific code generation at Layer 4. We briefly characterize the four
primary targets.

\subsection{Mesh Processors (GreenArrays-style)}

The GreenArrays GA144 contains 144 F18A cores in a mesh, each with 64 words of RAM,
an 18-bit instruction word, a data stack, and a return stack. Communication is
synchronous: a read from a port blocks until a write arrives from the neighboring node,
and vice versa. This blocking synchrony is the chip's primary coordination mechanism.

For this class, the cartographer extracts the mesh topology as a grid graph, assigns
Forth/colorForth instruction sets to each node, and models communication as synchronous
rendezvous with energy cost proportional to the number of transitions on the shared
port wire. The code generator emits colorForth-style source for each node.

\subsection{FPGAs}

FPGAs present a different constraint geometry: the computation itself is spatially
configurable, but routing resources are finite and congestion in the interconnect
fabric is a primary bottleneck. The cartographer models the FPGA fabric as a graph
of configurable logic blocks, DSP slices, block RAMs, and routing switches with
congestion-dependent latency and energy costs. The code generator emits HDL or a
high-level synthesis specification.

\subsection{GPUs}

GPUs offer massive parallelism but impose strict requirements on memory access patterns:
coalesced global memory access, shared memory bank-conflict avoidance, and occupancy
optimization. The cartographer models the GPU as a hierarchical graph of streaming
multiprocessors, each with shared memory and a warp-parallel execution model. The
program decomposer must respect the SIMT execution model. The code generator emits
CUDA or OpenCL kernels.

\subsection{Neuromorphic Chips}

Neuromorphic chips (Intel Loihi, IBM TrueNorth, SpiNNaker) compute through
spike propagation on fixed or configurable connection graphs. Energy is consumed
primarily at spike events, making sparsity of activation the central optimization
target. The cartographer models the neural network graph with synaptic weights and
routing tables. The code generator emits spike-routing configurations and
neuron parameter settings.

% -----------------------------------------------------------
\section{The Ecological Metaphor}

Moore's language suggests an ecological rather than mechanical understanding of the
chip: nodes occupy niches, communication channels are migration paths, energy is a
shared resource to be conserved, and the programmer's job is not to command execution
but to design an ecosystem in which the right computations emerge from local interactions.

This metaphor is not merely rhetorical. It suggests a generalization of the Silicon
Cartography framework toward \emph{machine ecology}: the study of how computational
load distributes across a spatial chip under realistic workloads, how communication
patterns create congestion or underutilization in particular chip regions, and how
programming choices shift the equilibrium of this distribution.

A machine ecology analysis of a compiled program would track:
\begin{itemize}
  \item \emph{Node utilization}: the fraction of time each node spends in each
        power state, revealing hot spots and idle regions;
  \item \emph{Edge saturation}: the fraction of available bandwidth consumed by
        each communication link, revealing routing bottlenecks;
  \item \emph{Entropy production}: the total energy expenditure per node per unit time,
        identifying which program components are thermodynamically expensive; and
  \item \emph{Dormancy patterns}: the temporal distribution of idle intervals,
        characterizing how well the program exploits the chip's low-energy states.
\end{itemize}

This analysis closes the loop with the hardware cartographer: the measured execution
ecology of a deployed program refines the cost model $P$, improving future compilations.

% -----------------------------------------------------------
\section{Formal Properties}

\subsection{Feasibility}

Not every program graph $G$ admits a valid compilation into every hardware graph $H$.
Feasibility requires that the kernel resource requirements $R$ are satisfiable within
the node capabilities $I$ and $M$, that data dependencies $D$ can be routed through
the edge set $E$ without violating bandwidth constraints, and that all timing constraints
in $C$ can be met. Feasibility checking is in general NP-complete (by reduction to
graph coloring), but for structured chip topologies---grids, trees, hypercubes---
polynomial-time algorithms exist.

\begin{theorem}[Feasibility on Grid Meshes]
For a program graph $G = (K, D, R)$ and a grid mesh hardware graph $H$ with uniform
node capabilities and edge capacities, feasibility is decidable in polynomial time
in $|K|$ and $|N|$ when all kernels require the same instruction types and
$|D| = O(|K|)$.
\end{theorem}

We omit the proof here; it follows from the fact that routing on a grid mesh with
bounded-degree program graphs reduces to a planar embedding problem.

\subsection{Optimality}

Finding $\varphi^*$ that minimizes $\text{Cost}(\varphi)$ over all valid compilations
is in general NP-hard (by reduction to quadratic assignment). In practice, for small
chips (fewer than a few hundred nodes), exact solutions via integer programming are
feasible. For larger chips, heuristic methods---simulated annealing, genetic algorithms,
message-passing belief propagation---find near-optimal placements efficiently.

\subsection{Compositionality}

A key property of the program graph formalism is compositionality: if $G_1$ and $G_2$
are program graphs with compatible message interfaces, their composition $G_1 \otimes G_2$
is a well-formed program graph. Valid compilations of $G_1$ and $G_2$ into disjoint
subsets of $N$ compose into a valid compilation of $G_1 \otimes G_2$, provided the
routing paths for inter-graph communication are feasible. This allows incremental
compilation and modular system design on spatial chips.

% -----------------------------------------------------------
\section{Fractal Topology Families}

\subsection{Beyond the Uniform Grid}

The GA144 is a flat rectangular mesh: eight rows by eighteen columns, with each F18A
node communicating through four blocking ports to its cardinal neighbors. This topology
is elegant and minimal, but it represents only one point in a much larger space of
spatial chip architectures. The Silicon Cartography framework is deliberately
topology-agnostic: the hardware graph $H$ accommodates any connected graph, and the
compilation functor adapts accordingly. This section develops a family of
procedurally generated chip topologies that extend the flat mesh toward recursive,
hierarchical, and fractal structures, each with characteristic locality properties
and programming disciplines.

\subsection{Fractal Island Mesh}

The simplest departure from the flat grid is recursive subdivision. Divide the
node budget into a $k \times k$ array of \emph{islands}, each containing a small
local mesh of worker nodes. Islands connect through dedicated \emph{bridge nodes}
whose sole function is inter-island routing. The resulting two-level hierarchy is
described by the generator

\[
N_{\text{island}} = N_{\text{local}}^{(k \times k)} \cup N_{\text{bridge}},
\]

where $N_{\text{bridge}}$ forms a secondary grid over the island topology. Bridge
nodes carry higher-bandwidth, higher-energy communication budgets than local worker
nodes, and may be assigned larger local memory to buffer inter-island traffic.

The programming discipline that this topology enforces is the central principle of
Silicon Cartography in recursive form: keep fast loops inside islands, export only
summaries across bridges, and never move raw state unless unavoidable. Algorithms
that exhibit local coherence and occasional long-range exchange---field relaxation,
graph diffusion, image processing, sparse neural inference---map naturally onto this
structure. The island boundary acts as an admissibility barrier, and the compiler's
locality preservation constraint

\[
d_G(\varphi(k_i), \varphi(k_j)) \leq L \cdot \|\sigma(k_i) - \sigma(k_j)\|
\]

is applied at two scales: within an island for local kernels and across islands for
summary reduction kernels.

The recursion may be applied again: a level-3 architecture organizes islands of
islands, with super-bridge nodes handling inter-cluster routing. This produces a
three-level hierarchy that mirrors the cache hierarchy of conventional processors
but with explicit programmer control over every level of the communication geometry.

\subsection{Tree-of-Meshes Architecture}

An alternative to recursive subdivision is the tree-of-meshes: a rooted tree in
which leaves perform sensing or local computation, branch nodes perform reduction
and routing, and the root node coordinates global state. Each leaf cluster contains
a small flat mesh; sibling clusters may carry optional lateral cross-links with
probability $p$, creating a small-world connectivity structure on top of the tree backbone.

The generator is:
\begin{enumerate}
  \item Instantiate one root node with elevated memory $M_{\text{root}} \gg M_{\text{leaf}}$.
  \item Attach $k$ branch nodes, each with intermediate memory $M_{\text{branch}}$.
  \item Attach $k$ leaf-cluster meshes to each branch node.
  \item With probability $p$, add lateral links between sibling clusters.
\end{enumerate}

This architecture is suited to reductions: many local readings aggregate into regional
summaries and ultimately into global decisions. Sensor swarms, audio analysis
pipelines, distributed event detection, and biological-style hierarchical inference all
exhibit this pattern. The tree structure makes the communication pattern maximally
explicit: every upward edge is a reduction, every downward edge is a control signal,
and lateral edges are optional coherence corrections.

The hardware graph for this topology has bounded degree at leaf nodes (four neighbors
in the local mesh plus one upward branch link) and higher degree at branch and root
nodes (up to $k$ downward links plus one upward link). The cost model must assign
appropriately higher energy budgets to branch and root nodes while enforcing that
leaf nodes spend the vast majority of their time dormant.

\subsection{Hilbert-Curve Chip}

A Hilbert curve generates a bijective map between the integers $\{0, \ldots, 2^{2n}-1\}$
and the cells of a $2^n \times 2^n$ grid that preserves locality: cells nearby in the
linear index are nearby in the grid. A Hilbert-curve chip is physically a flat grid,
but the compiler treats node addresses as Hilbert-curve indices, so that streaming
access along the logical order corresponds to physically local communication.

The Hilbert ordering is generated recursively. For a $4 \times 4$ grid the ordering
visits cells in the U-shaped pattern characteristic of the space-filling curve, and
at each recursive subdivision the pattern rotates and reflects to maintain the
locality property. The result is that any contiguous segment of the linear traversal
corresponds to a roughly square subregion of the grid, giving optimal locality for
both streaming and tiled computation.

For the compiler, this means that data with one-dimensional serial structure---
image scanlines, audio sample streams, compressed symbol sequences---should be
mapped onto the chip along the Hilbert index order. Operations that require spatial
locality (field relaxation, image tile convolution, compression context modeling)
automatically achieve physical locality as a consequence. The programmer specifies
the semantic ordering of data; the compiler uses the Hilbert bijection to find the
corresponding physical placement.

\subsection{Sierpinski Compute Fabric}

The Sierpinski-fabric architecture deliberately removes portions of the grid at
multiple scales. Given a base square lattice, the generator recursively excises
the center subregion at each level of subdivision, leaving a fractal structure
of connected compute cells separated by voids. These voids are not defects; they
serve as routing reservoirs, thermal gaps, analog interface corridors, or dedicated
memory channels.

At level $m$ the node set is:
\[
S_{m+1} = \bigcup_{i=1}^{3} f_i(S_m), \qquad f_i(x) = \tfrac{1}{2}x + v_i,
\]
where $v_1, v_2, v_3$ are the vertices of the base triangle and each $f_i$ is a
contraction by factor $\tfrac{1}{2}$. The resulting communication graph has a
characteristic fractal dimension $d_f = \log 3 / \log 2 \approx 1.585$ and
a degree distribution that is highly non-uniform: nodes near the boundary of each
sub-triangle have higher degree than interior nodes.

The power-density motivation is central. On conventional dense grids, active
regions near each other produce overlapping thermal fields that raise the local
temperature, increasing leakage current and degrading reliability. The Sierpinski
voids provide recursive thermal separation between active subregions, allowing
higher per-node power budgets in each compute cell without exceeding global thermal
limits. Programming for this topology concentrates dense computation within the
connected compute cells and uses boundary nodes as the only communication points
between sub-regions.

\subsection{Dragon-Curve Pipeline}

The dragon-curve architecture is a folded asynchronous pipeline. A long linear
computation---a filter chain, a signal processing cascade, an event detection
hierarchy---is physically folded into a compact recursive curve that fits within
the chip's spatial bounds. The folding preserves the logical order of the pipeline
while minimizing the physical wire length between adjacent pipeline stages.

The dragon curve is generated by repeatedly replacing each directed segment with
a right-angled pair of segments, alternating orientation at each level. The result
is a space-filling path with the property that each segment is adjacent to the
next in the logical pipeline ordering and also physically close to two or three
other pipeline stages. This locality enables a ``pass-forward'' programming discipline:
each node wakes when a value arrives from its predecessor, performs a local transform,
forwards the result to its successor, and returns to dormancy. Feedback is used only
at the recursive fold points, where the curve doubles back on itself.

This architecture is ideal for streaming computations where the data flow is
fundamentally one-directional but the chip is two-dimensional: audio sample
processing, image scanline filtering, sensor event pipelines, and data compression.

\subsection{Apollonian Heterogeneous Fabric}

The Apollonian architecture abandons node uniformity entirely. Drawing on the
geometry of circle packing, it places large high-capability nodes at the center of
the chip and recursively smaller specialized nodes around them. The largest central
region handles memory-intensive or globally coordinating computation. Medium regions
handle DSP, routing, and analog interfacing. Small peripheral nodes handle preprocessing
and event filtering. Tiny always-on nodes at the outermost periphery perform
wake-up detection at near-zero power.

The programming rule is: use tiny nodes for event detection, use small nodes for
preprocessing, use medium nodes to aggregate and route, and wake large nodes only
when the accumulated evidence makes full processing necessary. This matches the
inference pattern of biological sensory systems, in which peripheral neurons
perform massive signal reduction before forwarding only salient events to
higher processing regions.

% -----------------------------------------------------------
\section{Communication Manifold and Spectral Structure}

\subsection{The Chip as a Metrized Graph}

The hardware graph $G = (N, E)$ is not merely topological. It is
\emph{metrized}: each edge carries a cost that reflects the physical properties of
the wire, buffer, or bus it represents. Define a weighted edge cost

\[
w(i,j) = \alpha \ell_{ij} + \beta \tau_{ij} + \gamma \epsilon_{ij},
\]

where $\ell_{ij}$ is the geometric wire length, $\tau_{ij}$ is the communication
latency, and $\epsilon_{ij}$ is the energy required for a single data transfer,
with $\alpha, \beta, \gamma \geq 0$ weighting coefficients. This induces a
communication metric over all admissible routing paths:

\[
d_G(i,j) = \inf_{\pi: i \leadsto j} \sum_{e \in \pi} w(e).
\]

The metric $d_G$ is not symmetric in general: a chip may have asymmetric buses,
directional links, or one-way DMA channels. For undirected substrates, $d_G$
is a proper metric and equips the node set with a metric space structure that
captures the energetic geography of communication.

\subsection{Graph Laplacian and Diffusion}

Define the graph Laplacian $L = D - A$ where $A$ is the weighted adjacency matrix
and $D$ is the degree matrix with $D_{ii} = \sum_j w(i,j)$. The Laplacian governs
how disturbances propagate across the chip. A message broadcast or synchronization
signal injected at node $n_0$ diffuses according to:

\[
\frac{\partial u}{\partial t} = -Lu, \qquad u(0) = \delta_{n_0}.
\]

The eigenvalues $0 = \lambda_0 \leq \lambda_1 \leq \cdots \leq \lambda_{|N|-1}$
of $L$ characterize the chip's communication geometry. The algebraic connectivity
$\lambda_1$ (Fiedler value) measures how well-connected the chip is: a small
$\lambda_1$ indicates a bottleneck or near-partition, while a large $\lambda_1$
indicates high global connectivity. For fractal architectures, the spectrum
of $L$ encodes the self-similar structure of the topology and can be computed
recursively from the spectra of smaller instances.

The spectral decomposition of the chip graph informs compilation directly. Kernels
that must communicate frequently should be placed on nodes that are spectrally
close (small eigenvector distance in the low-frequency subspace). Dormant regions
correspond to parts of the chip that decouple from the active subgraph, appearing
as near-zero entries in the relevant eigenvectors.

\subsection{Synchronization Sheaves and Global Sections}

For asynchronous chips, communication events are not clock-driven but
handshake-driven. A port read on node $i$ blocks until a port write arrives from
node $j$. This synchronization constraint can be formalized using sheaf theory.

Define a presheaf $\mathscr{F}$ over the open sets of the chip graph $G$,
assigning to each subgraph $U \subseteq G$ the local state space
$\mathscr{F}(U) = \prod_{i \in U} S_i$ of all nodes in $U$. For overlapping
subgraphs $U$ and $V$ with non-empty intersection, the restriction maps require
consistency on the shared boundary:

\[
s_U \big|_{U \cap V} = s_V \big|_{U \cap V}.
\]

A valid distributed computation is then a global section $s \in \Gamma(G, \mathscr{F})$:
an assignment of local states to every node that is consistent across all
communication interfaces simultaneously. Deadlock corresponds to the
absence of a global section: some combination of blocking port states admits
no consistent resolution. The sheaf-theoretic formulation makes this failure mode
precise and suggests detection algorithms based on Čech cohomology obstruction
classes in $H^1(G, \mathscr{F})$.

\subsection{Iterated Function Systems and Recursive Chip Generation}

For fractal chip topologies, the node set itself is generated by an iterated
function system $\mathcal{F} = \{f_1, \ldots, f_k\}$ acting on an embedding
space $X \subseteq \mathbb{R}^2$. Each $f_i$ is a contraction map. The node
set at recursion depth $m$ is:

\[
N_m = \bigcup_{i_1, \ldots, i_m} f_{i_1} \circ \cdots \circ f_{i_m}(N_0).
\]

The communication graph at depth $m$ is the union of the images of the depth-$(m-1)$
graph under each $f_i$, together with the inter-image bridge edges that connect
adjacent subregions. This recursive structure means that the hardware graph $H_m$
satisfies a self-similarity relation:

\[
\mathbf{Chip}_m \simeq \coprod_{i=1}^{k} \mathbf{Chip}_{m-1}^{(i)},
\]

where the coproduct is in the category whose objects are spatial chip regions and
whose morphisms are admissible communication transformations. Compilation for
fractal chips exploits this recursive structure: a kernel decomposition for depth
$m$ is obtained by recursively solving the depth-$(m-1)$ problem for each
sub-region and then solving the inter-region routing problem for the bridge nodes.

% -----------------------------------------------------------
\section{Procedural Topology Language}

\subsection{A Unified Descriptor Schema}

The taxonomy of chip topologies developed in Section~5 suggests a general
\emph{procedural topology language}: a declarative schema in which a chip
architecture is specified not by enumerating nodes and edges but by giving
a generator, a recursion depth, a node-type assignment rule, and a cost model.
The compiler then instantiates the hardware graph $H$ from this specification
before beginning placement.

A chip specification in this language has five components. The \emph{topology
generator} is a rule for constructing the node set and edge set: flat grid,
recursive island, Hilbert curve, Sierpinski subdivision, dragon-curve pipeline,
or a composition thereof. The \emph{node capability assignment} maps node
positions in the generated topology to instruction sets, memory capacities, and
power domains, potentially varying by position (as in the Apollonian fabric where
node size decreases with distance from the center). The \emph{communication
geometry} specifies the edge types: blocking synchronous ports, buffered
asynchronous channels, shared-memory regions, or broadcast buses, each with
its own cost function $w(i,j)$. The \emph{energy model} assigns active and
idle power to each node type and transition type. The \emph{I/O boundary}
identifies which nodes interface with external pins, analog inputs, clock
sources, or off-chip memory.

Together these five components fully determine the hardware graph $H$ and its
cost model $P$, from which the cartographer layer constructs the executable
constraint map.

\subsection{Compiler Policy Parameters}

Beyond the hardware specification, the procedural topology language includes
a set of \emph{compiler policy parameters} that configure the optimization
behavior of the placement and routing layer:

\begin{itemize}
  \item \texttt{minimize\_instruction\_count}: prefer placements that require
        fewer total instructions by exploiting hardware-specific primitives;
  \item \texttt{minimize\_message\_distance}: place communicating kernel pairs
        on neighboring or nearby nodes;
  \item \texttt{prefer\_sleep\_over\_polling}: replace active waiting loops with
        blocking port reads wherever the communication pattern permits;
  \item \texttt{preserve\_locality}: enforce the Lipschitz condition
        $d_G(\varphi(k_i), \varphi(k_j)) \leq L \|\sigma(k_i) - \sigma(k_j)\|$
        with a specified Lipschitz constant $L$; and
  \item \texttt{thermal\_separation}: apply a minimum distance constraint
        $d_G(\varphi(k_i), \varphi(k_j)) \geq d_{\min}$ for kernel pairs
        whose simultaneous execution would exceed the local thermal budget.
\end{itemize}

These parameters allow domain-specific tuning. A sensor swarm application might
prioritize sleep maximization and message-distance minimization. A hard real-time
control system might override sleep preferences to guarantee bounded latency.
A thermally constrained wearable device might impose strict thermal separation
constraints. The procedural topology language thus separates hardware description
from compilation policy, making both independently reusable.

% -----------------------------------------------------------
\section{Visualization Framework}

\subsection{The Inadequacy of Block Diagrams}

A spatial chip in the sense defined here is not a collection of processors connected
by labeled arrows. It is a dynamic energetic manifold with locality structure,
recursive topology, wake/sleep state transitions, communication flow, and
admissibility constraints that vary continuously across the substrate. A conventional
block diagram is structurally inadequate to represent these six interacting
dimensions simultaneously:

\[
(\text{topology},\ \text{activity},\ \text{energy},\ \text{communication},\ \text{memory},\ \text{time}).
\]

A visualization framework suited to spatial chips must represent all six dimensions,
ideally at multiple spatial and temporal scales. This section develops such a framework.

\subsection{Spatial Graph Embedding with Physical Metrics}

The foundational layer is a spatial graph embedding $\varphi : N \to \mathbb{R}^2$
that places each node at its physical silicon location rather than at an
arbitrary graph-layout position. Node size encodes memory capacity:

\[
r_i \propto \log(1 + m_i),
\]

so that memory-rich nodes (branch and root nodes in a tree-of-meshes, central nodes
in an Apollonian fabric) are visually prominent. Edge thickness encodes bandwidth:

\[
w_{ij} \propto b_{ij},
\]

so that high-bandwidth bridges appear as thick channels and low-bandwidth
point-to-point links appear as thin filaments. The resulting image is not an
abstract graph diagram but a topographic map of the chip's resource geography.

\subsection{Activity Heatmaps and Dormancy Fields}

Activity visualization extends the spatial embedding with a temporal dimension.
Define the time-averaged activity of node $i$ over a computation trace of duration $T$:

\[
\bar{\chi}_i = \frac{1}{T} \int_0^T \chi_i(t)\, dt,
\]

where $\chi_i(t) \in \{0,1\}$ is the activity indicator. Rendering $\bar{\chi}_i$
as a color field over the spatial embedding produces an \emph{activity heatmap}:
high-activity nodes appear in warm colors, dormant nodes appear dark. A critical
insight for these architectures is that darkness is success. An activity heatmap
for a well-compiled program on a GreenArrays-like chip should show most nodes as
dark, with brief bright pulses propagating along the communication paths defined
by the program's data flow.

Thermal accumulation introduces a related field:

\[
T(x, t) = \int_0^t k_{\text{diff}}(x, x') \cdot \rho_e(x', t')\, dx'\, dt',
\]

where $\rho_e$ is the local energy dissipation density and $k_{\text{diff}}$
is a thermal diffusion kernel. Rendering $T(x,t)$ as a slowly evolving overlay
on the activity map reveals thermal hotspots that may not be visible in the
instantaneous activity alone.

\subsection{Event Wave Propagation}

For asynchronous chips, clock-cycle diagrams are meaningless: there is no global
clock. Instead, the appropriate temporal visualization is the \emph{event wave}:
only actual state transitions are rendered. A transition at node $i$ at time $t$
appears as a brief luminous pulse at position $\varphi(i)$. The pulse decays
with a time constant $\tau_{\text{vis}}$ chosen to make individual events visible
without creating persistent glow. When the computation involves a message cascade
(one node's output triggering a neighbor's computation triggering the next
neighbor's), the resulting visualization shows a wave of pulses propagating
across the chip surface, resembling neural spike propagation, mycelial signaling,
or fluid turbulence.

Repeated high-traffic communication paths accumulate visible luminosity over many
computation cycles, forming what might be called \emph{semantic highways}: the
chip's dominant information arteries become geometrically visible without any
explicit routing annotation.

\subsection{Causal Braid Diagrams}

For programs with complex synchronization structure, the event-wave visualization
obscures the causal order of transitions. A \emph{causal braid diagram} makes
causality explicit by assigning each node a vertical strand in a spacetime diagram
(time increasing downward) and representing communication events as crossings
between strands. The braid structure reveals synchronization dependencies,
bottlenecks, deadlock cycles, burstiness, and idle gaps in a form that is
directly analogous to string diagrams in monoidal category theory.

For asynchronous systems, causal braid diagrams are vastly more informative than
clock-cycle diagrams, because they show the actual causal dependency structure
of the computation rather than its alignment to an external time reference.
A computation with no unnecessary synchronization appears as a braid with
no unnecessary crossings; redundant synchronization points appear as spurious
entanglements that a compiler optimization pass can remove.

\subsection{Multiscale Semantic Zooming for Fractal Topologies}

For fractal chip topologies, a static rendering at any single scale is inadequate:
the recursive structure that defines the architecture's properties appears only
across multiple scales simultaneously. A \emph{multiscale semantic zoom} renders
the chip at multiple levels of detail simultaneously, with the visible level of
detail varying continuously as the viewer zooms in or out.

At the coarsest zoom level, entire islands or subtrees appear as single
representative nodes, with inter-island communication edges rendered as thick
channels. At intermediate zoom, the internal structure of each island becomes
visible, with local communication patterns emerging. At the finest zoom level,
individual stack states and instruction sequences appear. This is directly
analogous to geographic information systems in which political boundaries,
cities, roads, and buildings emerge at successive zoom levels.

For fractal topologies generated by iterated function systems, the self-similarity
of the structure means that the same visual grammar applies at every zoom level:
the viewer exploring a Sierpinski-fabric chip sees the same triangular-void
pattern at the scale of the whole chip, at the scale of a sub-region, and at
the scale of an individual compute cell's neighborhood.

\subsection{Energy Landscape Terrain}

Perhaps the most powerful single visualization for constraint-aware compilation
is the \emph{energy landscape terrain}: a three-dimensional surface over the
chip's spatial embedding in which the height $z(x,y,t)$ equals the instantaneous
energy dissipation rate $E(x,y,t)$. Rendering this surface as a dynamic terrain
map transforms the chip into a living landscape: active computation appears as
mountains and peaks, dormant regions appear as valleys and plains, and
communication events appear as traveling ridges.

Efficient programs produce sparse ripples and localized transient peaks that
quickly relax back to the valley floor. Inefficient programs produce flat global
elevation (all nodes consuming energy simultaneously), persistent plateaus
(nodes stuck in polling loops), or chaotic turbulence (synchronization failures
causing repeated retries). The terrain metaphor makes the energy cost of
programming decisions immediately visible and viscerally comprehensible.

This is a direct geometric expression of Moore's core principle: unnecessary
switching is waste, and waste is visible as unnecessary height in the energy
landscape. A programmer who can see the terrain their program creates on a
spatial chip can reason about energy in the same intuitive way that a civil
engineer reasons about drainage and runoff.

\subsection{Compiler Placement Geometry Overlay}

A final visualization layer represents the compiler's output directly: the
mapping $\varphi : K \to N$ from computational kernels to physical nodes.
Each kernel is rendered as a colored territorial overlay on the chip's spatial
embedding, with the kernel's color encoding its computational role (sensing,
reduction, output, control). Data dependencies between kernels appear as
elastic bands stretching between their territorial boundaries.

Compilation quality is then visually immediate. A placement with good locality
shows each kernel's territory as a compact, convex region, with short elastic
bands connecting communicating kernels. A placement with poor locality shows
kernels stretched across the chip, with long elastic bands crossing many
intermediate nodes. Thermal hotspots appear as multiple high-energy kernels
occupying adjacent territories. Synchronization bottlenecks appear as many
elastic bands converging on a single bridge node.

This overlay closes the loop between the cartographer's hardware model and the
programmer's algorithmic intent: it shows, in physical terms, exactly where the
computation lives and how information moves between its components.

% -----------------------------------------------------------
\section{Spacetime Activity and the Thermodynamic Cost of Computation}

\subsection{Activity Volume}

A spatially distributed computation induces a \emph{spacetime activity volume}:
the integral over time of the number of active nodes,

\[
\mathcal{V} = \int_0^T |A(t)|\, dt,
\]

where $A(t) \subseteq N$ is the set of nodes transitioning at time $t$. Minimizing
$\mathcal{V}$ is equivalent to minimizing total dynamic energy expenditure when
per-node transition costs are uniform. For non-uniform costs, the weighted
activity volume

\[
\mathcal{V}_w = \int_0^T \sum_{i \in A(t)} \eta_i(t)\, dt
\]

is the appropriate objective, where $\eta_i(t)$ is the instantaneous transition
energy of node $i$. This is the quantity that the cost functional $J(\varphi)$
approximates in the static compilation setting.

The activity volume formulation makes explicit what mainstream software metrics
obscure: the relevant measure of a computation's cost is not its wall-clock
duration (which conflates execution time with idle time) and not its
instruction count (which does not account for spatial parallelism) but the
total spacetime mass of its physical transitions.

\subsection{Sparse Activation as a Design Criterion}

The proposition established in Section~4---that energy minimality is equivalent
to trajectory sparsity---extends naturally to the spacetime framework.
A computation is \emph{maximally sparse} when at each instant $t$ the active
node set $A(t)$ is as small as possible while still making progress toward
the desired output. This corresponds to the ideal of asynchronous event-driven
computation: no node executes unless a communication event has made its
execution necessary.

Achieving maximal sparsity requires that the program graph $G$ be decomposed
into kernels with minimal data dependencies, that the compilation mapping
$\varphi$ places communicating kernels on neighboring nodes (minimizing
the number of intermediate routing hops that must activate), and that
the code generation layer replaces all polling loops with blocking port reads.
These three requirements---minimal dependency graph, locality-preserving placement,
blocking communication---are the core conditions of the Silicon Cartography discipline.

\subsection{Connection to Thermodynamic Computation Theory}

The identification of computation with physical state transitions connects
Silicon Cartography to the theoretical foundations of thermodynamic computation.
Landauer's principle establishes that erasing one bit of information requires
a minimum energy dissipation of $k_B T \ln 2$, where $k_B$ is Boltzmann's
constant and $T$ is the physical temperature. This bound applies to any
logically irreversible operation: any computation that discards information
must pay a thermodynamic price.

For practical computing on spatial chips, Landauer's bound is not yet the
limiting factor---switching energy in current CMOS technology exceeds the
Landauer limit by several orders of magnitude. But the direction of approach
is clear: as transistors shrink and switching energies fall, the thermodynamic
cost of logically irreversible operations will become increasingly significant.
Programming disciplines that minimize unnecessary state transitions---like the
one Moore describes and Silicon Cartography formalizes---are not merely
practical optimizations for the present. They are preparations for a
computational regime in which thermodynamic cost is not an engineering
approximation but the fundamental constraint.

In this sense, Moore's moral language about waste is not rhetorical. Every
unnecessary transistor transition is a physical irreversibility---a small but
real contribution to the entropy of the universe. A programmer who takes this
seriously is not performing premature optimization; they are acknowledging the
physical reality of computation.

% -----------------------------------------------------------
\section{Discussion}

\subsection{Recovery of Physical Meaning}

The central claim of Silicon Cartography is that energy-minimal computing on spatial
chips recovers a form of \emph{physical meaning} in programs that mainstream software
engineering has erased. In a conventional compiled language, the connection between
a line of source code and the physical transistor transitions it induces is opaque:
it passes through an optimizing compiler, an instruction scheduler, a microarchitectural
pipeline, a cache hierarchy, and a memory controller, each of which transforms the
computation in ways invisible to the programmer.

On a spatial chip programmed through Silicon Cartography, the connection is transparent.
Each kernel is physically located at a node. Each communication is a physical wire
transition. Each idle period is a physical dormancy. The cost model $P$ is not
an approximation of something deeper; it is the direct physical reality of the computation.

This transparency is not merely aesthetically satisfying. It is practically necessary
for programming chips at the energy frontier. A programmer who cannot reason about
the physical cost of their design choices cannot make them well.

\subsection{From Sequential Control to Constraint Propagation}

Traditional programming languages model computation as the ordered execution of
symbolic instructions over centralized mutable state. This perspective emerged from
serial machine architectures in which global synchronization was technologically
convenient, and it has been inherited by nearly every layer of the mainstream software
stack---operating systems, garbage collectors, memory models, sequential consistency
protocols. The question that rarely gets asked is whether this inheritance reflects
something necessary about computation or something contingent about a particular
engineering era.

The framework proposed here suggests the latter. Computation, viewed physically,
is constrained propagation across admissibility manifolds: programs are not sequences
of instructions but structured activation geometries embedded within hardware topology.
Under this interpretation, control flow becomes activation flow, memory locality
becomes geometric proximity, synchronization becomes compatibility of propagation
fronts, and scheduling becomes manifold navigation under energetic constraints.

The distinction is not merely semantic. Sequential models externalize coordination
into scheduling mechanisms, synchronization primitives, cache coherence protocols,
and global execution semantics. Constraint-propagation models internalize coordination
into the admissibility structure itself. This internalization eliminates an entire
class of coordination overhead that the sequential model must introduce and then
laboriously manage. The resulting architecture more closely resembles naturally
distributed systems---vascular networks, neural tissue, reaction-diffusion
systems---where coherent global behavior emerges through local propagation rules
and energetic admissibility conditions rather than centralized orchestration.

This shift has a historical parallel in physics: the transition from force-centric
to field-centric formulations. Just as modern physics replaced action-at-a-distance
mechanisms with continuous field descriptions, future computational systems may
increasingly replace imperative control structures with distributed admissibility
geometries governing local activation propagation. The field is not weaker than the
force; it is a more faithful account of what is actually happening.

An important strand of theoretical computer science has argued for a related
reconception from the symbolic side: that concurrency is simpler and more primitive
than sequentiality, and that sequential execution is best understood as a particular
serialization imposed over a fundamentally distributed substrate. Process-network
semantics, communicating-process algebras, dataflow models, and invocation-based
process expression all move in this direction at the symbolic level. The present
framework extends that reconception into explicitly physical territory: spatial
embedding, thermodynamic cost, energy topology, routing metrics, and hardware
admissibility constraints transform the philosophical argument into an engineering
discipline with concrete implementation consequences.

\subsection{The Thermodynamic Cost of Serialization}

Sequential execution imposes a thermodynamic burden that abstract computational
models routinely conceal. Serialization forces computational regions that could
otherwise evolve independently to synchronize around a globally ordered execution
schedule. This synchronization introduces energetic overhead through clock
distribution, pipeline stalling, cache coherence maintenance, global arbitration,
and repeated state serialization and deserialization. In sparse activation systems
operating near the energy frontier, these coordination costs can exceed the energy
spent on useful computation by orders of magnitude.

The present framework treats serialization not as a neutral abstraction but as a
physically consequential geometric restriction. A globally ordered execution
schedule constrains the dimensionality of allowable activation configurations:
at each moment, only the single active node specified by the program counter may
transition, while all other nodes remain artificially frozen. This is an extreme
restriction of the accessible manifold of low-energy computational states.

By contrast, asynchronous sparse activation permits computational activity to
distribute across all locally admissible regions simultaneously. Energy expenditure
concentrates near active computational frontiers while dormant regions remain
quiescent. The spacetime activity volume $\mathcal{V}_w$ is minimized not by
accelerating a single execution thread but by allowing the computation's natural
distributed geometry to evolve without artificial serialization constraints.

This observation suggests that many contemporary inefficiencies in large-scale
embedded and distributed computation may not arise primarily from transistor
limitations or memory bandwidth---the usual suspects---but from geometric mismatches
between sequential symbolic abstractions and the distributed thermodynamic structure
of physical hardware. Addressing these mismatches requires not faster serialization
but the dissolution of unnecessary serialization: a return to programming disciplines
that negotiate with the physical geometry of computation rather than abstracting it away.

\subsection{Three Senses of Locality}

The compiler policy parameter \texttt{preserve\_locality} encodes a Lipschitz
condition $d_G(\varphi(k_i), \varphi(k_j)) \leq L \|\sigma(k_i) - \sigma(k_j)\|$
that constrains the mapping from semantic to physical distance. This constraint
conceals an important conceptual distinction among three senses of locality that
the framework simultaneously manages.

\emph{Physical locality} is the most direct: two nodes are physically local when
they are adjacent or nearby in the hardware graph $G$, meaning that communication
between them costs little energy and imposes minimal latency. Physical locality is
determined entirely by chip geometry and is independent of the computation being performed.

\emph{Logical locality} is a property of the program graph $G$: two kernels are
logically local when they share a direct data dependency---when one kernel's output
is the immediate input of the other. Logical locality is determined by the structure
of the computation and is independent of the hardware.

\emph{Semantic locality} is the subtler quantity: two kernels are semantically local
when they operate on conceptually related data or perform functionally coherent
operations, even if they are not directly connected in the dependency graph.
A sensor kernel and its associated calibration kernel are semantically local even if
a reduction kernel sits between them in the dependency chain.

The Lipschitz condition above maps semantic distance $\|\sigma(k_i) - \sigma(k_j)\|$
to physical distance $d_G$, requiring that the placement preserve semantic geometry.
This is the deepest of the three locality conditions because it connects algorithmic
intent to physical substrate: a well-compiled program is one in which the chip's
physical geometry reflects the semantic structure of the computation it implements.

This connection has broad resonance. In sparse representation learning, the goal is
to find embeddings in which semantic similarity maps to geometric proximity in a
low-dimensional space. In biological cortical organization, functionally related
neurons are spatially clustered, and this clustering reduces the wiring length and
energy cost of correlated activations. In graph partitioning, the goal is to assign
computationally related tasks to the same processor to minimize inter-processor
communication. Silicon Cartography unifies these concerns under a single geometric
discipline: the placement $\varphi$ is an embedding of the program's semantic
geometry into the chip's physical geometry, and compilation quality is measured
by how faithfully this embedding preserves the relevant metric structure.

\subsection{Causal Geometry and Temporal Admissibility}

Time has been implicit throughout this framework in traces, schedules, and activity
volumes. For asynchronous spatial chips, where no global clock coordinates execution,
a more explicit treatment of the temporal structure of computation is necessary. The
appropriate formalism is a causal partial order.

\begin{definition}[Causal Precedence]
Given a computation trace on hardware graph $H$, define the \emph{causal precedence}
relation $\prec$ on execution events $(k_i, t_i)$ by:

\[
(k_i, t_i) \prec (k_j, t_j)
\]

if and only if kernel $k_j$ depends on output from $k_i$ and the information
emitted by $k_i$ at time $t_i$ can physically reach $\varphi(k_j)$ before
time $t_j$, i.e.\ $t_j \geq t_i + d_G(\varphi(k_i), \varphi(k_j)) / c_{\max}$,
where $c_{\max}$ is the maximum signal propagation speed on the chip.
\end{definition}

The causal precedence relation defines a partial order on execution events: it is
irreflexive, asymmetric, and transitive. The forward causal cone of an event
$(k_i, t_i)$ is the set of all events that can be influenced by it:

\[
\mathcal{C}^+(k_i, t_i) = \{(k_j, t_j) \mid (k_i, t_i) \prec (k_j, t_j)\}.
\]

The causal cone is determined jointly by the dependency structure of $G$ and the
communication metric $d_G$ of $H$. Two events are \emph{causally disjoint} when
neither lies in the other's causal cone; causally disjoint events may execute
simultaneously without coordination overhead, because no information can flow
between them within the relevant time interval.

This formulation makes the hardware graph into something closer to a finite
asynchronous spacetime manifold. Physical distance in $d_G$ bounds the speed of
causal propagation; the chip's topology determines which computations can run
concurrently and which must be sequenced. A placement $\varphi$ that maximizes the
causal disjointness of computationally independent kernels minimizes unnecessary
serialization and therefore minimizes the thermodynamic cost of coordination.

Deadlock can now be characterized geometrically: a deadlock occurs when a set of
events $\{(k_i, t_i)\}$ form a cycle in the causal precedence relation, so that
each event is waiting for a predecessor that is itself waiting. This is equivalent
to the non-existence of a global section in the synchronization sheaf
$\mathscr{F}$---consistent with the Čech cohomology obstruction interpretation
developed in Section~6. The causal and sheaf-theoretic analyses thus
converge on the same necessary condition for deadlock-free computation.

Temporal admissibility adds a further constraint to the compilation problem.
A valid compilation must not only satisfy spatial admissibility (kernel $k$
at node $n$ with $R(k) \subseteq I(n)$) but also temporal admissibility:
the causal precedence constraints must be satisfiable within the timing windows
specified in $C$. Bounded-latency proofs for safety-critical systems can then
be stated as geometric assertions about the causal cone structure of the
compiled program: a kernel $k_j$ meets its deadline if and only if all
predecessors in $\mathcal{C}^-(k_j, t_j^{\max})$ can reach $\varphi(k_j)$
by time $t_j^{\max}$.

\subsection{A Worked Miniature Example}

To anchor the preceding abstractions, consider a concrete miniature instantiation
of the framework on a nine-node island mesh.

\paragraph{Hardware.} The hardware graph $H_9$ consists of nine F18A-style nodes
arranged in a $3 \times 3$ grid. Interior node $n_5$ (center) connects to all four
cardinal neighbors. Corner nodes ($n_1, n_3, n_7, n_9$) each have two neighbors.
Edge nodes ($n_2, n_4, n_6, n_8$) each have three neighbors. Node $n_1$ serves as
the external input port; node $n_9$ serves as the external output port.
Each edge has uniform energy cost $w = 1$ and latency $\tau = 1$. Each node has
64 words of local memory and supports stack-machine instruction execution.

\paragraph{Program.} The task is a sensor-filter-reduce-output pipeline:
a scalar sensor reading enters at $n_1$, passes through a digital filter (a
five-tap FIR computation requiring five multiply-accumulate operations), is reduced
by a threshold comparison, and the binary result is emitted at $n_9$.

The program graph $G$ has four kernels: $k_{\text{sense}}$ (sample and forward
the input), $k_{\text{filter}}$ (accumulate five taps), $k_{\text{reduce}}$
(threshold comparison), and $k_{\text{out}}$ (format and emit output). The
dependency chain is linear: $k_{\text{sense}} \to k_{\text{filter}} \to
k_{\text{reduce}} \to k_{\text{out}}$.

\paragraph{Valid placement.} The locality-preserving placement assigns:
$\varphi(k_{\text{sense}}) = n_1$, $\varphi(k_{\text{filter}}) = n_2$,
$\varphi(k_{\text{reduce}}) = n_3$, $\varphi(k_{\text{out}}) = n_6$.
Each communication hop traverses one edge; total communication cost is $3w = 3$.
Nodes $n_4, n_5, n_7, n_8, n_9$ remain dormant throughout. The activity volume
for a single pipeline invocation is $\mathcal{V} = 4$ (four kernel executions) plus
$3$ (three intermediate communication events) $= 7$ node-time units.

\paragraph{Invalid placement.} An alternative placement assigns
$\varphi(k_{\text{filter}}) = n_7$, routing the sensor output from $n_1$ to $n_7$
through $n_4$ and $n_7$ (two hops, passing through $n_4$) and the filtered output
back from $n_7$ to $n_3$ via $n_4$ or $n_8$ (two more hops). The intermediate
routing nodes $n_4$ must activate to relay messages even though they perform no
computation. Total communication cost rises to $5w = 5$; the activity volume
rises to approximately $11$ node-time units due to relay activations. This
placement also violates the thermal-separation constraint if $k_{\text{filter}}$
and $k_{\text{sense}}$ were assigned to the same voltage domain boundary.

\paragraph{Energy comparison.} The valid placement achieves an activity volume
of $7$ units; the invalid placement achieves approximately $11$ units---a $57\%$
increase in energy expenditure for identical computational output. On a chip
running at $10^9$ pipeline invocations per second, this difference corresponds to
a sustained power increase of several milliwatts, which is significant at the
energy budget of a sensor swarm node. The activity heatmap for the valid placement
shows five dark (dormant) nodes and four briefly active nodes arranged along the
top edge of the grid; the invalid placement shows six active nodes including two
relay activations whose sole contribution is message forwarding.

\subsection{Relation to Existing Frameworks}

Silicon Cartography does not emerge from a vacuum. Several prior and parallel
research traditions have addressed aspects of the problem it addresses, and
positioning the framework relative to these traditions is necessary to identify
what is genuinely new.

Dataflow programming languages---from Dennis's original data flow procedure
language through Kahn process networks and Lucid---decompose computation into
stream-processing networks in which nodes activate when input data is available.
This is structurally similar to the kernel-decomposition model, and the blocking
communication discipline of GA144-style chips is directly implementable as a
Kahn network. The difference is that Silicon Cartography adds explicit physical
embedding: the placement of dataflow nodes onto chip geometry with energy-cost
optimization is not addressed by dataflow semantics, which abstract away the
physical substrate entirely.

Communicating sequential processes and related process algebras (CCS, $\pi$-calculus)
provide compositional models of concurrent communication that closely parallel the
synchronization sheaf formulation. The handshake admissibility condition $r_i(t)
\wedge w_j(t) = 1$ is a direct instance of CSP synchronization. The difference
again is physicality: process algebras operate on abstract named channels; Silicon
Cartography operates on metrized hardware edges with energy, latency, and congestion.

Systolic arrays formalize the mapping of regular array computations onto spatially
regular processor meshes, achieving high throughput through careful data scheduling
and local communication. The systolic principle---that data should flow rhythmically
through a regular array without requiring global control---is closely related to
the streaming computation model of Hilbert-curve and dragon-curve chips. The
difference is that systolic arrays require globally synchronous clocking, while
Silicon Cartography targets asynchronous chips and prioritizes energy minimality
over throughput maximization.

Network-on-chip research addresses the design and optimization of the on-chip
communication fabric for many-core processors, including topology selection,
routing algorithms, and arbitration policies. The hardware graph formalism directly
subsumes the network-on-chip problem as a special case: the chip's communication
topology is modeled as a metrized graph with bandwidth, latency, and contention
constraints. The difference is that network-on-chip research typically assumes a
fixed program structure and optimizes the network for it, while Silicon Cartography
co-optimizes program decomposition and network placement simultaneously.

High-level synthesis tools (Halide, TVM, XLA) optimize the mapping of array
computations onto hardware targets including GPUs, FPGAs, and custom accelerators.
These tools perform schedule and placement optimization within their respective
computational models. Silicon Cartography generalizes beyond regular array
computation to arbitrary program graphs, handles asynchronous communication
explicitly, and treats energy as the primary optimization target rather than
throughput or latency.

Place-and-route tools in VLSI design solve the problem of assigning logic gates
to physical locations and routing wires between them---a problem that is structurally
similar to the compilation mapping $\varphi : G \to H$. The key difference is level
of abstraction: place-and-route operates at the level of logic gates and wires;
Silicon Cartography operates at the level of computational kernels and communication
protocols. The formalisms are compatible and a Silicon Cartography compiler could
in principle generate a place-and-route specification for an FPGA or custom ASIC.

The genuinely novel synthesis offered by Silicon Cartography is the unification of
these concerns under a single geometric-thermodynamic ontology: spatial compilation,
thermodynamic cost modeling, admissibility geometry, fractal topology generation,
synchronization sheaves, causal partial orders, and visualization as energetic
cartography. No existing framework addresses all of these simultaneously, and the
mutual reinforcement among them---particularly the way that fractal topology
generation, spectral analysis, sheaf-theoretic deadlock detection, and energy
landscape visualization all operate on the same underlying hardware graph $H$---is
the framework's most distinctive structural contribution.

\subsection{Limitations and Open Problems}

The framework as presented assumes that the hardware graph $H$ is fully known at
compile time. For reconfigurable chips (FPGAs, some neuromorphic systems), $H$ may
itself be a design variable, opening a joint optimization problem over chip
configuration and program placement. Extending Silicon Cartography to this setting
requires a richer formalism in which the compilation functor maps into a space of
(hardware configuration, program placement) pairs.

A second open problem is the handling of dynamic workloads: inputs that vary
in size, shape, or arrival rate over time. The static compilation model produces
a fixed placement; adapting it dynamically requires either pre-compiling a family
of placements indexed by workload parameters, or developing online replanning
algorithms that can migrate kernels across nodes at runtime.

A third open problem concerns the procedural topology language itself: given a
computational workload and an energy budget, what chip topology should be
generated? The inverse design problem---find $\mathcal{F}$ and recursion depth $m$
such that the induced hardware graph $H_m$ admits a low-cost compilation of a
given program class $\mathcal{G}$---is structurally a program synthesis problem
over a topology grammar, and is not addressed by any of the existing tools
reviewed above.

Finally, the causal geometry framework raises the question of reversible computation.
Fredkin and Toffoli demonstrated that logically reversible computation can in
principle approach the Landauer limit of zero dissipation per operation. Integrating
reversible computation primitives into the Silicon Cartography kernel model---so
that kernels can be annotated as reversible where the computation permits---would
allow the framework to reason about the thermodynamic floor of a given computation,
not merely its current energy cost.

% -----------------------------------------------------------
\section{Conclusion}

Silicon Cartography proposes a discipline inversion for programming spatial chips:
begin with the machine, not the program. By modeling the hardware as a metrized
constraint graph, decomposing tasks into locally self-contained kernels with
explicit communication interfaces, and finding energy-minimal placements through
constrained optimization over the chip's admissibility field, it recovers the
direct negotiation between algorithm and matter that abstraction-heavy software
engineering has systematically obscured.

The framework operates simultaneously at three intellectual levels that mutually
constrain one another. At the practical level it provides a four-layer compiler
architecture with concrete algorithms for hardware cartography, capability modeling,
program decomposition, and placement routing. At the mathematical level it furnishes
a unified set of tools---metrized hardware graphs, graph Laplacians, synchronization
sheaves, iterated function systems for fractal topology generation, causal partial
orders, and spacetime activity integrals---that make each architectural decision
formally precise. At the philosophical level it argues that the dominant tradition
of sequential symbolic computation is not a neutral description of what computation
is, but a historically contingent engineering choice whose thermodynamic costs
become increasingly visible as computing moves toward the energy frontier.

The fractal topology families developed here---island meshes, tree-of-meshes,
Hilbert-curve chips, Sierpinski compute fabrics, dragon-curve pipelines, Apollonian
fabrics---demonstrate that the hardware graph formalism is not limited to the
flat uniform grid of the GA144 but encompasses a wide design space of spatial
substrates, each with characteristic spectral properties, locality structures,
and programming disciplines. The procedural topology language provides a common
descriptor schema for this design space, separating hardware geometry from
compilation policy and enabling systematic exploration of chip-program co-design.

The visualization framework translates the abstract mathematical structure of
spatial computation into directly perceivable form: activity heatmaps make
dormancy visible, event wave propagation makes causal structure visible, energy
landscape terrain makes thermodynamic cost visible, and causal braid diagrams
make synchronization structure visible. Together these representations constitute
a cartographic discipline in the literal sense: they map the invisible energetic
geography of a running computation onto a form that can be directly navigated.

The connection to thermodynamic computation theory---Landauer's principle,
conservative logic, the activity volume integral as a measure of irreversibility
expenditure---grounds the energy-minimality criterion in physics rather than
engineering preference. Moore's moral language about waste is, in the end,
thermodynamically precise: every unnecessary transistor transition is a small
irreversibility, a minor but real entropic cost paid by the physical universe
for the privilege of discarding information. A programming discipline that takes
this seriously is not performing premature optimization. It is computing honestly.

% -----------------------------------------------------------
\section*{Acknowledgments}

The author thanks the extended discussions around GreenArrays documentation,
Forth philosophy, and constraint-first computing that informed this framework.

% -----------------------------------------------------------
\begin{thebibliography}{99}

\bibitem{moore2013}
C.\ Moore, ``GreenArrays Architecture and colorForth Programming,''
unpublished lecture notes and documentation, GreenArrays Inc., 2013.

\bibitem{moore1974}
C.\ Moore, ``Programming a Problem-Oriented Language,''
unpublished manuscript, 1974.

\bibitem{greenarr}
GreenArrays Inc., ``GA144 Product Brief,'' 2011.
\url{http://www.greenarraychips.com}

\bibitem{fant2005}
K.\ M.\ Fant, \textit{Computer Science Reconsidered: The Invocation Model of
Process Expression}, Wiley-IEEE Computer Society Press, 2005.

\bibitem{fant1985}
K.\ M.\ Fant and S.\ A.\ Brandt, ``Null Convention Logic: A Complete and Consistent
Logic for Asynchronous Digital Circuit Synthesis,'' in \textit{Proceedings of the
International Conference on Application Specific Array Processors}, 1985.

\bibitem{martin1990}
A.\ J.\ Martin, ``The Limitations to Delay-Insensitivity in Asynchronous Circuits,''
in \textit{Advanced Research in VLSI}, MIT Press, 1990, pp.\ 263--278.

\bibitem{sparso2001}
J.\ Spars{\o} and S.\ Furber, \textit{Principles of Asynchronous Circuit Design:
A Systems Perspective}, Kluwer Academic Publishers, 2001.

\bibitem{dennis1974}
J.\ B.\ Dennis, ``First Version of a Data Flow Procedure Language,''
in \textit{Programming Symposium}, Lecture Notes in Computer Science, vol.\ 19,
Springer, 1974, pp.\ 362--376.

\bibitem{wadge1985}
W.\ W.\ Wadge and E.\ A.\ Ashcroft, \textit{Lucid, the Dataflow Programming Language},
Academic Press, 1985.

\bibitem{kahn1974}
G.\ Kahn, ``The Semantics of a Simple Language for Parallel Programming,''
in \textit{Proceedings of IFIP Congress 74}, North-Holland, 1974, pp.\ 471--475.

\bibitem{hoare1978}
C.\ A.\ R.\ Hoare, ``Communicating Sequential Processes,''
\textit{Communications of the ACM}, 21(8):666--677, 1978.

\bibitem{milner1989}
R.\ Milner, \textit{Communication and Concurrency}, Prentice Hall, 1989.

\bibitem{backus1978}
J.\ Backus, ``Can Programming Be Liberated from the von Neumann Style?
A Functional Style and Its Algebra of Programs,''
\textit{Communications of the ACM}, 21(8):613--641, 1978.

\bibitem{landauer1961}
R.\ Landauer, ``Irreversibility and Heat Generation in the Computing Process,''
\textit{IBM Journal of Research and Development}, 5(3):183--191, 1961.

\bibitem{bennett1982}
C.\ H.\ Bennett, ``The Thermodynamics of Computation---A Review,''
\textit{International Journal of Theoretical Physics}, 21(12):905--940, 1982.

\bibitem{fredkin1982}
E.\ Fredkin and T.\ Toffoli, ``Conservative Logic,''
\textit{International Journal of Theoretical Physics}, 21(3--4):219--253, 1982.

\bibitem{mead1980}
C.\ Mead and L.\ Conway, \textit{Introduction to VLSI Systems},
Addison-Wesley, 1980.

\bibitem{hennessy2019}
J.\ L.\ Hennessy and D.\ A.\ Patterson, \textit{Computer Architecture: A Quantitative
Approach}, 6th ed., Morgan Kaufmann, 2019.

\bibitem{dally2001}
W.\ J.\ Dally and B.\ Towles, ``Route Packets, Not Wires: On-Chip Interconnection
Networks,'' in \textit{Proceedings of the 38th Design Automation Conference},
ACM/IEEE, 2001, pp.\ 684--689.

\bibitem{benini2002}
L.\ Benini and G.\ De Micheli, ``Networks on Chips: A New SoC Paradigm,''
\textit{Computer}, 35(1):70--78, 2002.

\bibitem{fiedler1973}
M.\ Fiedler, ``Algebraic Connectivity of Graphs,''
\textit{Czechoslovak Mathematical Journal}, 23(2):298--305, 1973.

\bibitem{chung1997}
F.\ R.\ K.\ Chung, \textit{Spectral Graph Theory},
American Mathematical Society, 1997.

\bibitem{barnsley1988}
M.\ F.\ Barnsley, \textit{Fractals Everywhere}, Academic Press, 1988.

\bibitem{sagan1994}
H.\ Sagan, \textit{Space-Filling Curves}, Springer, 1994.

\bibitem{hilbert1891}
D.\ Hilbert, ``{\"U}ber die stetige Abbildung einer Linie auf ein
Fl{\"a}chenst{\"u}ck,'' \textit{Mathematische Annalen}, 38(3):459--460, 1891.

\bibitem{watts1998}
D.\ J.\ Watts and S.\ H.\ Strogatz, ``Collective Dynamics of Small-World Networks,''
\textit{Nature}, 393:440--442, 1998.

\bibitem{spivak2014}
D.\ I.\ Spivak, \textit{Category Theory for the Sciences}, MIT Press, 2014.

\bibitem{maclane1998}
S.\ Mac Lane, \textit{Categories for the Working Mathematician},
2nd ed., Springer, 1998.

\bibitem{ghrist2014}
R.\ Ghrist, \textit{Elementary Applied Topology}, Createspace, 2014.

\bibitem{curry2012}
J.\ Curry, R.\ Ghrist, and M.\ Robinson, ``Euler Calculus with Applications to
Signals and Sensing,'' \textit{Proceedings of Symposia in Applied Mathematics},
70:75--145, 2012.

\bibitem{robinson2014}
M.\ Robinson, \textit{Topological Signal Processing}, Springer, 2014.

\bibitem{merolla2014}
P.\ A.\ Merolla et al., ``A Million Spiking-Neuron Integrated Circuit with a
Scalable Communication Network and Interface,''
\textit{Science}, 345(6197):668--673, 2014.

\bibitem{davies2018}
M.\ Davies et al., ``Loihi: A Neuromorphic Manycore Processor with On-Chip
Learning,'' \textit{IEEE Micro}, 38(1):82--99, 2018.

\bibitem{furber2014}
S.\ Furber et al., ``The SpiNNaker Project,''
\textit{Proceedings of the IEEE}, 102(5):652--665, 2014.

\bibitem{kogge1994}
P.\ M.\ Kogge et al., ``Processing in Memory: Chips to Petaflops,''
Workshop on Mixing Logic and DRAM, 1994.

\bibitem{leiserson1991}
C.\ E.\ Leiserson and J.\ B.\ Saxe, ``Retiming Synchronous Circuitry,''
\textit{Algorithmica}, 6(1):5--35, 1991.

\bibitem{juarrero1999}
A.\ Juarrero, \textit{Dynamics in Action: Intentional Behavior as a Complex System},
MIT Press, 1999.

\bibitem{garey1979}
M.\ R.\ Garey and D.\ S.\ Johnson, \textit{Computers and Intractability: A Guide
to the Theory of NP-Completeness}, W.\ H.\ Freeman, 1979.

\end{thebibliography}

\end{document}
